\documentclass[aspectratio=169,9pt]{beamer}
\mode<presentation> 
\usepackage[utf8]{inputenc}
\usepackage{sansmathaccent} %Für equations
\usepackage{amsmath}
\usepackage{tikz}
\usepackage{multicol}
\usepackage{cancel}
\usepackage{relsize}
\usepackage{graphicx}
\usepackage[table]{xcolor}
\usepackage{fontawesome5}
\usepackage{enumitem}
\usepackage{setspace}
\usepackage[labelfont={bf},font={color=black!40,footnotesize},tablename={Tabelle},figurename={Abbildung}]{caption}
\usetikzlibrary{positioning, arrows.meta, calc, matrix}
\usefonttheme[onlymath]{serif}
\usepackage{booktabs}
%\makeatletter
%\makeatother

\usetikzlibrary{patterns, arrows.meta, calc}
\usepackage{pgfplots}
\pgfplotsset{compat=1.18}
\usetikzlibrary{intersections, pgfplots.fillbetween}\usetikzlibrary{intersections, pgfplots.fillbetween}
\rowcolors{2}{miugray!30}{miugray!10}

\usetheme{Madrid}  %% Themenwahl
\useoutertheme{split} % Alternatively: miniframes, infolines, split
%\useinnertheme{circles}


\setbeamertemplate{itemize item}{\color{miugray}$\blacksquare$}
\setbeamertemplate{itemize subitem}{\color{uniblau}$\bullet$}



%\newcommand{\metermilli}{\mskip3mu \text{mm}}
\newcommand{\cm}{\mskip3mu \text{cm}}
\newcommand{\m}{ \text{m}}
\newcommand{\lite}{ \text{l}}
\newcommand{\kg}{ \text{kg}}
\newcommand{\g}{ \text{g}}
\newcommand{\km}{ \text{km}}
\newcommand{\h}{ \text{h}}
\newcommand{\s}{ \text{s}}
\newcommand{\tonne}{ \text{t}}


\newcommand{\RA}{$\rightarrow$~}
\newcommand{\pp}{\par \vspace*{2mm}}
\newcommand{\ppbig}{\par \vspace*{4mm}}

\definecolor{miudiutex}{HTML}{1d2926}
\definecolor{miugray}{HTML}{b4cccf}
\definecolor{miugray2}{HTML}{4d7365}
\definecolor{white2}{HTML}{f5f7f8}
\definecolor{red}{HTML}{cf1f5c}
\definecolor{green}{HTML}{00A36C}
\definecolor{delphinidin}{HTML}{315794}
\definecolor{uniblau}{HTML}{004e9f}
\definecolor{unigelb}{HTML}{fcba00}
\definecolor{unigrau}{HTML}{909085}
\definecolor{pelargonidin}{HTML}{b33807}
\definecolor{cyanidin}{HTML}{ba0746}
\definecolor{paeonidin}{HTML}{d15ca6}
\definecolor{petunidin}{HTML}{b56fed}
\definecolor{malvidin}{HTML}{8a2d8c}

%\setbeamercolor{palette primary}{bg=pelargonidin!35,fg=white}
%\setbeamercolor{palette secondary}{bg=delphinidin!45,fg=white}
\setbeamercolor{palette primary}{bg=uniblau!70,fg=white}
\setbeamercolor{palette secondary}{bg=uniblau!45,fg=white}
\setbeamercolor{palette tertiary}{bg=unigelb!70,fg=unigrau}
\setbeamercolor{palette quaternary}{bg=miugray,fg=white}
\setbeamercolor{structure}{fg=black} % itemize, enumerate, etc
\setbeamercolor{section in toc}{fg=red} % TOC sections

\setbeamersize{text margin left=1cm,text margin right=1cm}


\useinnertheme{tcolorbox} 
\newcommand{\goodtk}[2]{
\begin{tcolorbox}[
	enhanced,
	colframe=uniblau!75,
	sharp corners,
	boxsep=5pt,
	title=\bfseries \sffamily #1,
	colback=white,
	coltitle=black,
	attach boxed title to top text left={yshift=-\tcboxedtitleheight/2},
	boxed title style={colback=white,colframe=white}
]
	#2
	\end{tcolorbox}
}
	
\newcommand{\antwort}[1]{
\begin{tcolorbox}[
	enhanced,
	colframe=unigrau!65,
	sharp corners,
	boxsep=1pt,
	colback=white,
	coltitle=black,
]
	#1
	\end{tcolorbox}
}


\title{11. Übung}
\author{PD Dr. Helene Loos}
\date{\today}

\begin{document}
%\maketitle
\small



\begin{frame}
\begin{center}

\LARGE Übung 11  \par \normalsize
Prozessbezogene Lebensmitteltechnologie
\end{center}
\small
\end{frame}



\begin{frame}


\begin{minipage}{.5\textwidth}
{\color{miugray}$\blacksquare ~$} \textbf{Aufgabe  1} \par \vspace*{2mm}

\begin{enumerate}[label=\alph*)]
\item<1-> Was versteht man unter dem pH-Wert? Wie ist er definiert
\item<3-> Stellen Sie die Gleichung für die Gleichgewichtskonstante der Reaktion einer Bronstedt-Säure HA mit Wasser auf.
\item<5-> Wie berechnet man den pH-Wert für eine wässrige Lösung einer starken Säure (pK$_s <$ 1)?
\item<7-> Wie berechnet man den pH-Wert für eine wässrige Lösung einer schwachen Säure?
\item<9-> Definieren Sie den Dissoziationsgrad einer Säure
\item<11-> Wie viel Prozent der Moleküle einer in Wasser gelösten Säure liegen dissoziiert vor, wenn der pK$_S$-Wert dem pH-Wert entspricht?
\end{enumerate}


\end{minipage}%
\hspace*{4mm}
\begin{minipage}{.5\textwidth}
\antwort{
\begin{enumerate}[label=\alph*)]
\item<2-> pH = $-\log(\text{H}^+)$ 
\item<4-> HA + H$_2$O $\leftrightharpoons$ A$^-$ + H$_3$O$^+$
\item<6-> Wir gehen davon aus, dass HA komplett dissoziert. Daher pH = $-\log(\text{H}^+)$
\item<8-> pH = $\frac{1}{2} (pK_S - \log(c_0))$
\item<10-> $\alpha = \frac{[A^-]}{[HA]}$
\item<12-> Henderson-Hasselbalch: \pp pH = $pK_S + \log \left(\frac{A^-}{HA} \right)$ \RA Also 50 \%
\end{enumerate}
}

\end{minipage}%
\end{frame}






\begin{frame}


\begin{minipage}{.5\textwidth}


{\color{miugray}$\blacksquare ~$} \textbf{Aufgabe  2} \par \vspace*{2mm}

Bei der Entwicklung einer Light Mayonnaise mit einem Fettgehalt von 35\% steht die Frage im Raum, ob Sorbinsäure oder Benzoesäure zugesetzt werden soll. Der pH-Wert der Mayonnaise liegt bei pH 4,4.
\begin{enumerate}[label=\alph*)]
\item[a)] Berechnen Sie, wie hoch jeweils der Anteil an dissoziierter Säure wäre. Der pKS-Wert
von Benzoesäure liegt bei 4,18 und der pK$_S$-Wert von Sorbinsäure bei 4,76.
\end{enumerate}
\end{minipage}%
\hspace*{4mm}
\begin{minipage}{.5\textwidth}

\antwort{

\uncover<2->{Säurekonstante nach Dissoziationsgrad $\alpha$ umstellen:}
\begin{align*}
\uncover<2->{K &= \frac{[H^+] \cdot [A^-]}{[HA]} \\}
\uncover<3->{K &= \frac{[H^+] \cdot \alpha}{1 - \alpha} \\}
\uncover<4->{K \cdot (1-\alpha) & = [H^+] \cdot \alpha \\}
\uncover<5->{K - \alpha \cdot K &= [H^+] \cdot \alpha \\}
\uncover<6->{K &= [H^+] \cdot \alpha + \alpha \cdot K \\}
\uncover<7->{K &= \alpha \cdot ([H^+] + K) \\}
\uncover<8->{\alpha &= \frac{K}{[H^+] + K}}
\end{align*}
}


\end{minipage}%






\end{frame}





\begin{frame}


\begin{minipage}{.5\textwidth}


{\color{miugray}$\blacksquare ~$} \textbf{Aufgabe  2} \par \vspace*{2mm}

Bei der Entwicklung einer Light Mayonnaise mit einem Fettgehalt von 35\% steht die Frage im Raum, ob Sorbinsäure oder Benzoesäure zugesetzt werden soll. Der pH-Wert der Mayonnaise liegt bei pH 4,4.
\begin{enumerate}[label=\alph*)]
\item[a)] Berechnen Sie, wie hoch jeweils der Anteil an dissoziierter Säure wäre. Der pKS-Wert
von Benzoesäure liegt bei 4,18 und der pK$_S$-Wert von Sorbinsäure bei 4,76.
\end{enumerate}
\end{minipage}%
\hspace*{4mm}
\begin{minipage}{.5\textwidth}

\antwort{

\begin{align*}
\alpha &= \frac{K}{[H^+] + K}
\end{align*}

Nun können wir einfach einsetzen mit $K= 10^{-pK_S}$ und $[H^+] = 10^{-pH}$ \pp
\uncover<2->{
\textbf{Benzoesäure:}
\begin{align*}
\alpha &= \frac{10^{-4,18}}{10^{-4,4} + 10^{-4,18}} \\
&= 0,624
\end{align*}
Daher liegt der Anteil an undissoziierter Säure bei 1-0,624 = 0,3759 \RA 37,6 \%
}
\pp 
\uncover<3->{
\textbf{Sorbinsäure:}
\begin{align*}
\alpha &= \frac{10^{-4,76}}{10^{-4,4} + 10^{-4,76}} \\
&= 0,303
\end{align*}
Daher liegt der Anteil an undissoziierter Säure bei 1-0,303 = 0,697 \RA 69,7 \%
}
}
\end{minipage}%

\end{frame}







\begin{frame}

\begin{minipage}{.5\textwidth}
{\color{miugray}$\blacksquare ~$} \textbf{Aufgabe  2} \par \vspace*{2mm}

Bei der Entwicklung einer Light Mayonnaise mit einem Fettgehalt von 35\% steht die Frage im Raum, ob Sorbinsäure oder Benzoesäure zugesetzt werden soll. Der pH-Wert der Mayonnaise liegt bei pH 4,4.
\begin{enumerate}[label=\alph*)]
\item[b)] Welche Relevanz hat das Ergebnis für die Auswahl des Konservierungsstoffs? Welche
weiteren Aspekte sind bei der Auswahl von Konservierungsstoffen zu beachten?
\end{enumerate}



\end{minipage}%
\hspace*{4mm}
\begin{minipage}{.5\textwidth}

\antwort{
\begin{itemize}
\item[\color{miugray}$\blacktriangleright$]<2-> Für konservierende Wirkung ist der Anteil an undissozierter Säure ausschlaggebend. Für den vorliegenden pH-Wert wäre also die Sorbinsäure stärker konservierend wirksam.
\item[\color{miugray}$\blacktriangleright$]<3-> Betrachtung der ADI-Werte / Höchstmengen
\item[\color{miugray}$\blacktriangleright$]<4-> Wirkungsspektrum und ''Ziel-MO''
\item[\color{miugray}$\blacktriangleright$]<5-> Möglichkeit des kombinierten Einsatzes
\item[\color{miugray}$\blacktriangleright$]<6-> Sensorische Eigenschaften des Produktes
\item[\color{miugray}$\blacktriangleright$]<7-> Verteilungskoeffizient Fett/Wasser
\end{itemize}

}

\end{minipage}%

\end{frame}






\begin{frame}

\begin{minipage}{.5\textwidth}
{\color{miugray}$\blacksquare ~$} \textbf{Aufgabe  2} \par \vspace*{2mm}
\begin{enumerate}[label=\alph*)]
\item[c)]  Nehmen Sie an, dass 200 mg/kg der Säure undissoziiert in der Mayonnaise vorliegen. Wie hoch ist der Anteil in der wässrigen Phase? Verwenden Sie zur Abschätzung den Verteilungskoeffizient K$_{OW}$ der beiden Säuren. Für Benzoesäure beträgt er in der
Mayonnaise 4,5, für Sorbinsäure beträgt er 8,7.
\end{enumerate}

\end{minipage}%
\hspace*{4mm}
\begin{minipage}{.5\textwidth}

\antwort{
\begin{align*}
\uncover<2->{K=4,5 &= \frac{\frac{200~\text{mg}-x}{350}}{\frac{x}{650~\text{g}}} \\}
\uncover<3->{4,5 &= \frac{200~\text{mg}-x}{350~\text{g}} \cdot \frac{650~\text{g}}{x}  \\}
\uncover<4->{4,5 &= \frac{650}{350} \cdot (\frac{200~\text{mg}}{x}-1) \\ }
\uncover<5->{4,5 \cdot \frac{350}{650} &= \frac{200~\text{mg}}{x} -1 \\}
\uncover<6->{(4,5 \cdot \frac{350}{650}+1) &= \frac{200~\text{mg}}{x} \\}
\uncover<7->{\frac{1}{4,5 \cdot \frac{350}{650} + 1} \cdot 200 ~\text{mg} &= x  \\}
\uncover<8->{x &= 58,43  ~\text{mg}}
\end{align*}
\uncover<9->{
Bei gleicher Verfahrensweise erhält man für \textbf{Sorbinsäure} (Verteilungskoeffizient = 8,7) einen Anteil von 35 mg in der wässrigen Phase}
}

\end{minipage}%


\end{frame}



\begin{frame}

\begin{minipage}{.5\textwidth}
{\color{miugray}$\blacksquare ~$} \textbf{Aufgabe  3} \par \vspace*{2mm}

Der pK$_S$-Wert von Buttersäure beträgt 4,82. Berechnen Sie den Anteil an undissoziierter Säure
bei den pH-Werten 1, 2, 3, 4, 5 und 6 und 7.
\end{minipage}%
\hspace*{4mm}
\begin{minipage}{.5\textwidth}
\antwort{
\footnotesize
\uncover<2->{
Henderson-Hasselbalch:
\begin{align*}
pK_S &= pH -\log(\frac{[A^-]}{HA})
\end{align*}
}
\uncover<3->{
Umstellen:
}
\begin{align*}
\uncover<4->{pH - pK_S &= \log(\frac{[A^-]}{HA}) \\}
\uncover<5->{10^{pH - pK_S} &= \frac{[A^-]}{HA} \\}
\uncover<6->{	&~=~ \frac{1 - [HA]}{HA} \\}
\uncover<7->{	&~=~ \frac{1}{HA} - \frac{[HA]}{HA} \\}
\uncover<8->{	&~=~ \frac{1}{HA} - 1 \\}
\\
\uncover<9->{10^{pH - pK_S}	&~=~ \frac{1}{HA} - 1 \\}
\uncover<10->{10^{pH - pK_S} +1 &= \frac{1}{HA} \\ \\}
\uncover<11->{HA &= \frac{1}{10^{pH - pK_S} + 1 }}
\end{align*}
}
\end{minipage}%
\end{frame}




\begin{frame}

\begin{minipage}{.5\textwidth}
{\color{miugray}$\blacksquare ~$} \textbf{Aufgabe  3} \par \vspace*{2mm}

Der pK$_S$-Wert von Buttersäure beträgt 4,82. Berechnen Sie den Anteil an undissoziierter Säure
bei den pH-Werten 1, 2, 3, 4, 5 und 6 und 7.
\end{minipage}%
\hspace*{4mm}
\begin{minipage}{.5\textwidth}
\antwort{
\uncover<2->{
\begin{align*}
HA &= \frac{1}{10^{pH - pK_S} + 1 }
\end{align*}
\pp
Nachdem Einsetzen von pK$_S$ und pH-Wert erhält man folgende Ergebnisse: \pp
pH 1 \RA 99,985\% \par pH 2 \RA 99,85\% \par pH 3 \RA 98,5\% \par pH 4 \RA 86,9\% \par pH 5 \RA 39,8\% \par pH 6 \RA 6,2\% \par pH 7 \RA 0,656\%
}
}
\end{minipage}%
\end{frame}



\begin{frame}

\begin{minipage}{.5\textwidth}
{\color{miugray}$\blacksquare ~$} \textbf{Aufgabe  4} \par \vspace*{2mm}

Propansäure hat einen pK$_S$-Wert von 4,88. Berechnen Sie den pH-Wert einer 0,1 molaren
Propansäurelösung.

\end{minipage}%
\hspace*{4mm}
\begin{minipage}{.5\textwidth}
\antwort{
Formel für schwache Säuren:
\begin{align*}
\uncover<2->{pH &= \frac{1}{2} \cdot (pK_S - \log(c)) \\}
\uncover<3->{pH &= \frac{1}{2} \cdot (4,88 - \log(0,1)) = 2,9}
\end{align*}


}
\end{minipage}%
\end{frame}




\begin{frame}
\begin{center}
\Large Herleitungen zu Henderson-Hasselbalch und schwache Säuren
\end{center}
\end{frame}

\begin{frame}
\begin{minipage}{.45\textwidth}
{\color{miugray}$\blacksquare ~$} \textbf{Titrationskurve } \par \vspace*{2mm}

\includegraphics[width=6cm]{../misc/titration.pdf}

\end{minipage}%
\hspace*{4mm}
\begin{minipage}{.55\textwidth}
\antwort{
\footnotesize
\begin{align*}
K_S &= \frac{c[H_3O^+] \cdot c[A^-]}{c[HA]} \\
K_S &= c[H_3O^+] \cdot \frac{c[A^-]}{c[HA]} \\
\log_{10}(K_S) &= \log_{10} (c[H_3O^+]) + \log_{10} \left(\frac{c[A^-]}{c[HA]}\right) \\
- \log_{10} (c[H_3O^+]) &= -\log_{10}(K_S) + \log_{10} \left(\frac{c[A^-]}{c[HA]}\right) \\
pH &= pK_S + \log_{10} \left(\frac{c[A^-]}{c[HA]}\right)
\end{align*}
}
\end{minipage}%



\end{frame}



\begin{frame}


\begin{minipage}{.5\textwidth}
{\color{miugray}$\blacksquare ~$} \textbf{Herleitung Formel Schwache Säuren  } \par \vspace*{2mm}

\antwort{
Eine schwache Säure reagiert in Wasser unter folgender Gleichung: 
\begin{align*}
HA + H_2O \leftrightharpoons H_3O^+ + A^-
\end{align*}
Daraus kann man über das Massenwirkungsgesetz folgende Beziehung für die Säurestärke aufstellen:
\begin{align*}
K_S &= \frac{c[H_3O^+] \cdot c[A^-]}{c[HA]} \\
\end{align*}
Vereinfacht:
\begin{align*}
K_S &= \frac{c[H_3O^+]^2}{c[HA]} \\
\end{align*}
}

\end{minipage}%
\hspace*{4mm}
\begin{minipage}{.5\textwidth}
\antwort{
Nun nimmt man an, dass bei schwachen Säuren die undissoziierte Säuremenge in etwa der Ausgangskonzentration entspricht. Es gilt $c[HA] \approx c_0[HA]$. Daraus kann man wie folgt ableiten:
\begin{align*}
c[H_3O^+]^2 &= K_S \cdot c_0[HA] \\
c[H_3O^+] &= \left( K_S \cdot c_0[HA]\right) ^{\frac{1}{2}} \\
\log(c[H_3O^+]) &= \frac{1}{2} \cdot \left( \log(K_S) + \log(c_0[HA]) \right) \\
- pH &= \frac{1}{2} \cdot (-pK_S + \log(c_0[HA]) )\\
pH &= \frac{1}{2} \cdot (pK_S - \log(c_0[HA]))
\end{align*}
}
\end{minipage}%

\end{frame}





\begin{frame}

\begin{center}
\LARGE
Vielen Dank fürs Zuhören!
\end{center}

\end{frame}







\end{document}
